\section*{\texorpdfstring{\S\CurrentExerciseGroup}{Section \CurrentExerciseGroup}}
\phantomsection
\addcontentsline{toc}{section}{\texorpdfstring{Exercises for \S\CurrentExerciseGroup}{Exercises for Section \CurrentExerciseGroup}}

\begin{enumerate}
\def\labelenumi{\arabic{enumi})}
\item
  Assume that the hypotheses of Th. 1 of No.~1 are satisfied. Let \(h\) be a locally \(\mu\) -integrable function such that \(p\) is \((h \cdot \mu)\) -proper. Show that the function \(b \mapsto g(b) = \int h(t) d\lambda_b(t)\) , defined locally almost everywhere in B (for the measure \(\nu = p(\mu)\) ) is such that \(p(h \cdot \mu) = g \cdot \nu\) .
\item
  Let B be a locally compact space, \((\nu_{\iota})_{\iota\in I}\) the family of all positive measures on B. Let T be the product space \(I\times B\) , I being equipped with the discrete topology, and let \(\nu\) be the measure on T whose restriction to \(\{\iota\}\times B\) is the canonical image of the measure \(\nu_{\iota}\) on B, for all \(\iota\in I\) . Let p be the projection of T onto B; show that if there exists an uncountable compact subset of B, then \(\nu\) does not admit a pseudo-image measure under p.~(Show that every point of B would have measure \textgreater0 for such a measure.)
\item
  Give an example of a positive measure \(\mu\) on a Polish locally compact space \(T\) and a continuous mapping \(p\) of \(T\) into a Polish locally compact space \(B\) , such that \(p\) is not \(\mu\) -proper.
\item
  Let \(\mathrm{T}\) be the interval {[}0,1{]} of \(\mathbf{R}\) , \(\mu\) the Lebesgue measure on \(\mathrm{T}\) . Let \(\mathrm{R}\{x,y\}\) be the equivalence relation \(x - y \in \mathbf{Q}\) in \(\mathrm{T}\) . Show that \(\mathrm{R}\) is not \(\mu\) -measurable (apply Th. 3 of No.~4), but the graph of \(\mathrm{R}\) in \(\mathrm{T} \times \mathrm{T}\) is negligible for the product measure \(\mu \otimes \mu\) .
\item
  Let T be the union of Cantor's triadic set \(K \subset [0,1]\) and the interval I = {[}1,2{]} of R, and let \(\mu\) be the measure induced on the compact space T by Lebesgue measure. Let P be a nonmeasurable subset of I having the power of the continuum (Ch. IV, §4, Exer. 8), and let \(\psi\) be a bijection of K onto P. Consider the equivalence relation R in T for which every point x not belonging to \(K \cup P\) is its own equivalence class, and the class of a point \(y \in K\) consists of y and \(\psi(y)\) . Show that R is \(\mu\) -measurable, but the saturation of K for R is not \(\mu\) -measurable.
\item
  Let T be a Polish locally compact space, \(\mu\) a positive measure on T, R a \(\mu\) -measurable equivalence relation in T, p a \(\mu\) -measurable mapping of T into a Polish locally compact space B such that \(p(x) = p(y)\) is equivalent to \(R\{x, y\}\) , \(\nu\) a pseudo-image measure of \(\mu\) under p, and \(b \mapsto \lambda_b\) a disintegration of \(\mu\) by R. For every \(b \in B\) such that \(\lambda_b \neq 0\) , let \(C(b)\) be the class mod R that carries \(\lambda_b\) ; show that \(\varphi_{C(b)} \cdot \mu\) is proportional to \(\lambda_b\) . Give an example where \(\varphi_{C(b)} \cdot \mu = 0\) for every \(b \in B\) .
\item
  Let T be a Polish locally compact space, \(\mu\) a bounded positive measure on T, and R a \(\mu\) -measurable equivalence relation in T. Let \(\Omega\) be the subset of the space \(\mathcal{M}(\mathrm{T})\) of measures on T, formed by the measures \(\lambda \geqslant 0\) , of total mass \(\leqslant 1\) and nonzero; \(\Omega\) is locally compact when it is equipped with the vague topology (Ch. III, §1, No.~9, Cor. 2 of Prop. 15). Show that there exists one and only one positive measure \(\rho\) on \(\Omega\) such that: \(1^{\circ} \mu = \int_{\Omega} \lambda d\rho(\lambda)\) ; \(2^{\circ} \rho\) is concentrated on a subset \(B_{0}\) of \(\Omega\) whose elements are measures of total mass 1, carried by the classes mod R, two distinct elements of \(B_{0}\) being carried by distinct classes. (Consider a disintegration \(b \mapsto \lambda_{b}\) of \(\mu\) by the relation R, such that all of the \(\lambda_{b}\) have total mass 1, and the image of B under the mapping \(b \mapsto \lambda_{b}\) of B into \(\Omega\) ; make use of Th. 4.)
\item
  Let T be the interval \([0,2]\) in R, \(\mu\) the Lebesgue measure on T. Let A be a non-Borel set contained in the Cantor set \(K \subset [0,1]\) (cf.~GT, IV, §2, No.~4, Example and IX, §6, Exer. 6). Let S be the equivalence relation in T whose equivalence classes are the sets \(\{x\}\) for \(x \notin A \cup (A + 1)\) and the sets \(\{x, x + 1\}\) for \(x \in A\) . Show that S is \(\mu\) -measurable but that there does not exist any Borel section for S.
\item
  Let \(\mathbf{K}\) be a metrizable compact space, \(f\) a continuous mapping of \(\mathbf{K}\) into a Hausdorff topological space \(\mathbf{E}\) , and \(k\) any integer \(> 1\) . Denote by \(\mathrm{B}_k\) the subset of \(\mathbf{E}\) formed by the \(y\) such that \(\stackrel{-1}{f(y)}\) contains at least \(k\) distinct points; show that \(\mathrm{A}_k = \stackrel{-1}{f(\mathrm{B}_k)}\) is a Borel set in \(\mathbf{K}\) (for every integer \(n > 0\) , let \(\mathrm{B}_{kn}\) be the subset of \(\mathbf{E}\) formed by the \(y\) such that \(\stackrel{-1}{f(y)}\) contains at least \(k\) points whose mutual distances are all \(\geqslant 1/n\) ; show that \(\mathrm{B}_{kn}\) is closed).
\end{enumerate}

¶ 10) Let K be a metrizable compact space, μ a positive measure on K, and f a μ-measurable mapping of K into a Hausdorff topological space E. Denote by I (resp. U) the subset of E formed by the y such that \(\stackrel{-1}{f(y)}\) is infinite (resp. uncountable).

\begin{enumerate}
\def\labelenumi{\alph{enumi})}
\item
  Show that there exists in K a \(\mu\) -measurable set H such that \(f(\mathrm{H}) = \mathrm{I}\) and such that, for every \(y \in \mathrm{E}\) , \(\stackrel{-1}{f(y)} \cap \mathbb{C}\mathrm{H}\) is finite. (Let \((\mathrm{K}_n)\) be an increasing sequence of compact subsets of K such that the restriction of \(f\) to each \(\mathrm{K}_n\) is continuous and the complement N of the union F of the \(\mathrm{K}_n\) is \(\mu\) -negligible; for every \(k > 1\) , let \(\mathrm{B}_k\) be the subset of E formed by the \(y\) such that \(\mathrm{F} \cap \stackrel{-1}{f(y)}\) contains at least \(k\) distinct points, and let \(\mathrm{A}_k = \mathrm{F} \cap \stackrel{-1}{f(\mathrm{B}_k)}\) ; take for H the union of the set \(\mathrm{A} = \bigcap_k \mathrm{A}_k\) and the set \(\mathrm{N} \cap \stackrel{-1}{f(\mathrm{I})}\) , and use Exer. 9.)
\item
  Let \(\mathfrak{F}\) be the set of \(\mu\) -measurable sets \(M \subset H\) such that \(M \cap f(y)\) is finite for every \(y \in E\) ; let \(\alpha\) be the supremum of the measures \(\mu(M)\) for \(M \in \mathfrak{F}\) ; there exists an increasing sequence \((\mathrm{M}_n)\) of sets in \(\mathfrak{F}\) such that \(\lim_{n\to \infty}\mu (\mathrm{M}_n) = \alpha\) . Let \(\mathrm{P} = \bigcup_{n}\mathrm{M}_{n}\)
\end{enumerate}

and let \(S\) be a measurable section of \(H \cap \mathbb{C}P\) for the equivalence relation \(f(x) = f(y)\) (Th. 3); show that \(\mu(S) = 0\) .

\begin{enumerate}
\def\labelenumi{\alph{enumi})}
\setcounter{enumi}{2}
\tightlist
\item
  Show that there exists a \(\mu\) -negligible set \(Z \subset K\) such that \(f(Z) = U\) and that, for every \(\varepsilon > 0\) , there exists a \(\mu\) -measurable set \(L \subset K\) such that \(\mu(L) \leqslant \varepsilon\) and \(f(L) = I\) (observe that \(f(H \cap \mathbb{C}M_n) = I\) for every \(n\) and that \(U \subset f(S)\) ).
\end{enumerate}

¶ 11) Let K be a metrizable compact space, μ a positive measure on K, f a μ-measurable mapping of K into a locally compact space T, and ν a positive measure on T; I and U have the same meaning as in Exer. 10.

\begin{enumerate}
\def\labelenumi{\alph{enumi})}
\item
  Suppose that for every \(\mu\) -negligible set \(N \subset K\) , \(f(N)\) is \(\nu\) -negligible. Then \(U\) is \(\nu\) -negligible and, for every \(\mu\) -measurable subset \(A \subset K\) , \(f(A)\) is \(\nu\) -measurable.
\item
  In order that I be \(\nu\) -negligible and the image under \(f\) of every \(\mu\) -negligible set be \(\nu\) -negligible, it is necessary and sufficient that \(f\) satisfy the following condition: for every \(\varepsilon > 0\) , there exists a \(\delta > 0\) such that, for every subset \(A \subset K\) that is \(\mu\) -measurable and is such that \(\mu(A) \leqslant \delta\) , \(f(A)\) is \(\nu\) -measurable and \(\nu(f(A)) \leqslant \varepsilon\) . (To see that the condition is necessary, argue by contradiction, by considering a decreasing sequence \((A_n)\) of \(\mu\) -measurable sets whose intersection is \(\mu\) -negligible, but such that \(\nu(f(A_n)) \geqslant \alpha > 0\) ; then take the intersection of the \(f(A_n)\) and \(\mathbb{C}\mathrm{I}\).)
\end{enumerate}

¶ 12) Let E be the non-metrizable compact space obtained by equipping the interval \([-1,+1]\) of R with the topology T defined in Exer. 13 a) of GT, IX, §2; let \(\varphi\) be the mapping \(x \mapsto |x|\) of E onto I = {[}0,1{]} (equipped with the topology induced by that of R).

\begin{enumerate}
\def\labelenumi{\alph{enumi})}
\item
  Show that \(\varphi\) is continuous and that, denoting by \(\mathfrak{F}\) the set of subsets of \(E\) of the form \(\overline{\varphi}(A)\) , where \(A\) is a Borel subset of \(I\) , the Borel subsets of \(E\) are the subsets \(B\) such that there exists an \(M \in \mathfrak{F}\) for which \(B \cap \mathbb{C}M\) and \(M \cap \mathbb{C}B\) are countable (observe that every open set of \(E\) belongs to \(\mathfrak{F}\) ).
\item
  Show that for a numerical function \(f\) defined on \(\mathbf{E}\) to be continuous (for \(\mathcal{T}\) ), it is necessary and sufficient that it be regulated and that \(f(x-) = f((-x)+)\) for every \(x \in \mathbf{E}\) . One therefore defines a positive measure \(\mu\) on \(\mathbf{E}\) by setting \(\int f d\mu = \int_0^1 f(x) dx\) . Show that the image of \(\mu\) under \(\varphi\) is the Lebesgue measure on \(\mathbf{I}\) , and the \(\mu\) -negligible sets are the subsets negligible for Lebesgue measure on \([-1, +1]\) .
\item
  Deduce from \(a)\) and \(b)\) that there does not exist a \(\mu\) -measurable section for the \(\mu\) -measurable equivalence relation \(\varphi(x) = \varphi(y)\) in \(E\) .
\end{enumerate}

